% ====================================================================
%  EXECUTIVE SUMMARY
% ====================================================================
\chapter*{Executive Summary}
\addcontentsline{toc}{chapter}{Executive Summary}
\label{ch:executive-summary}

% ====================================================================
%  §1  INTRODUCTION
% ====================================================================
\section*{Introduction}

Deep neural networks (DNNs) drive a growing range of applications, yet
their computational cost places stringent demands on both throughput
and energy efficiency.
Dedicated spatial architectures (SAs)---custom accelerators built
around arrays of processing elements (PEs)---exploit data-level
parallelism to meet these demands.
The performance of an SA, however, depends critically on the
\emph{mapping}: a schedule that determines how loop iterations are
tiled, ordered, and distributed across PEs and memory levels.
For a single convolutional layer on a moderately complex SA, the
resulting \emph{map space} already exceeds $10^{24}$--$10^{32}$
candidates, motivating analytical cost models and automated
mappers~\cite{parashar2019timeloop, Factorflow}.

Conventional frameworks, such as Timeloop~\cite{parashar2019timeloop}
and FactorFlow~\cite{Factorflow}, operate on a \emph{single layer at
a time}: every intermediate activation tensor is written to off-chip
DRAM after one layer completes and read back when the next begins.
Because DRAM accesses are roughly two orders of magnitude more
expensive than on-chip SRAM accesses in terms of energy, this
ping-pong pattern dominates the total energy budget for many
workloads.
\emph{Layer fusion} eliminates this bottleneck by computing multiple
consecutive layers in a single, interleaved schedule, keeping
intermediate activations entirely on-chip.
Hardware implementations such as DepFiN~\cite{goetschalckx2023depfin}
demonstrate that depth-first, fused execution can substantially
reduce DRAM traffic.

However, no prior tool extends the full generality of the FactorFlow
factor-manipulation paradigm to multi-layer fusion.
This thesis fills that gap: FactorFlow is extended with
\emph{aliasing}, \emph{decoupling}, a \emph{depth-first heuristic}, and a
\emph{per-layer cost model} to represent, prune, and evaluate the
fused-layer map space.
The extended tool is validated against the DepFiN accelerator and
exercised on a systematic design-space exploration covering two
architectures, four CNN workloads, and three fusion modes.


% ====================================================================
%  §2  PROBLEM FORMULATION
% ====================================================================
\section*{Problem Formulation}

\paragraph{Why fusion changes the problem.}
A single convolution layer has $|D| = 6$ loop dimensions
($M, C, P, Q, R, S$).
When $F$ layers are fused into a unified loop nest, the dimension
count grows as:
%
\begin{equation}
  |D_{\text{fused}}| \;=\;
  \underbrace{4F}_{\substack{\text{per-layer:}\\
    C_i, R_i, S_i, Z_i}} \;+\;
  \underbrace{2(F-1)}_{\substack{\text{intermediate}\\
    \text{spatial: } Y_i, X_i}} \;+\;
  \underbrace{2}_{\substack{\text{output}\\
    P, Q}}
  \;=\; 6F.
  \tag{Dimension growth}
\end{equation}
%
For the 11-layer FSRCNN fusion used in validation, this yields
$|D_{\text{fused}}| = 66$ --- an order-of-magnitude increase over 6.

\paragraph{Map-space explosion.}
The map-space cardinality is dominated by:
%
\[
  |\mathcal{MS}| \;=\;
  \underbrace{(|D|!)^{L_m}}_{\text{loop permutations}}
  \;\cdot\;
  \underbrace{
    \prod_{d \in D}\;\prod_{p \mid d}
      \binom{v_p(d) + L - 1}{L - 1}
  }_{\text{factor allocations}},
\]
%
where $L_m$ is the number of memory levels and $L$ the total levels
(memory + spatial).
Table~\ref{tab:es-mapspace} quantifies the growth.

\begin{table}[htbp]
  \centering
  \caption{Map-space comparison: single layer vs.\ 11-layer fusion
           (DepFiN-like architecture).}
  \label{tab:es-mapspace}
  \small
  \begin{tabular}{l r r}
    \toprule
    \textbf{Quantity}
      & \textbf{Single layer}
      & \textbf{11-layer fusion} \\
    \midrule
    Dimensions ($|D|$)          & 6     & 66 \\
    Memory levels ($L_m$)       & 4     & 4 \\
    Total levels ($L$)          & 5     & 26 \\
    Permutation term            & $\sim\!10^{11}$
                                & $\sim\!10^{370}$ \\
    Factor-allocation term      & $\sim\!10^{13}$--$10^{20}$
                                & $>\!10^{93}$ \\
    \addlinespace
    \textbf{Total map-space}    & $\sim\!10^{24}$--$10^{32}$
                                & $>\!10^{463}$ \\
    \bottomrule
  \end{tabular}
\end{table}

\noindent
The fused map space exceeds $10^{463}$ --- over 400 orders of
magnitude larger than the single-layer case, making unguided search
entirely intractable.

\paragraph{Inter-layer tile-size dependencies.}
Fusion also introduces cascading constraints.
An output tile of size $P \times Q$ for the last layer forces each
preceding layer~$i$ to produce a progressively larger tile:
%
\begin{equation}
  \text{Layer } i: \quad
  \bigl(P + \textstyle\sum_{j=i+1}^{F-1}(R_j - 1)\bigr)
  \;\times\;
  \bigl(Q + \textstyle\sum_{j=i+1}^{F-1}(S_j - 1)\bigr),
  \tag{Cascading tiles}
\end{equation}
%
so deeper fusion tightens on-chip memory constraints for every layer.

\paragraph{Five operand categories.}
Unlike single-layer mapping (input, output, weights), fused mapping
distinguishes five operand types: (1)~global input, (2)~per-layer
weights $W_0, \dots, W_{F-1}$, (3)~intermediate output $A_i$,
(4)~intermediate input $A_{i-1}$, and (5)~global output.
Each memory level must specify bypass rules for all five categories;
the bypass configuration determines whether the depth-first fusion
guarantee (no intermediate DRAM traffic) is satisfied.


% ====================================================================
%  §3  RELATED WORK
% ====================================================================
\section*{Related Work}

\textbf{Layer-by-layer frameworks.}\;
Timeloop~\cite{parashar2019timeloop} introduced the mapper--model
paradigm for spatial architectures, and
FactorFlow~\cite{Factorflow} reformulated map-space exploration
via factor manipulation and the extended Einsum
notation~\cite{een}, enabling linear, quadratic, and hybrid
mappers.
Both operate on a single layer at a time.

\textbf{Fused-layer frameworks.}\;
DeFiNES~\cite{DeFiNES} extends ZigZag~\cite{zigzag} with
depth-first scheduling and inter-layer data-reuse analysis but
constrains the set of admissible dataflows.
LoopTree~\cite{LoopTree} generalises Timeloop to multi-level,
multi-operation loop nests but does not incorporate the
factor-based map-space navigation of FactorFlow.
Neither tool offers the aliasing/decoupling abstraction developed
here.

\textbf{Hardware implementations.}\;
DepFiN~\cite{goetschalckx2023depfin} is a 22\,nm depth-first
CNN accelerator with a $16 \times 128$ weight-stationary PE
array, a 1056\,KB Feature Memory (FMEM), and a 524\,KB Weight
Memory (WMEM), executing fused convolutional pipelines on-chip.

\textbf{Graph-level scheduling.}\;
Optimus~\cite{Optimus} addresses the complementary problem of
\emph{which} layers to fuse, but relies on coarse energy
estimates and does not optimise the dataflow mapping within each
fused segment.


% ====================================================================
%  §4  PROPOSED SOLUTION
% ====================================================================
\section*{Proposed Solution}

The extension rests on three modelling concepts and a generalised
cost model, all integrated into the existing FactorFlow codebase.

\paragraph{1.\ Aliasing.}
Each intermediate activation $A_i$ appears twice in the unified loop
nest: as the \texttt{int\_out} of layer~$i$ and the \texttt{int\_in}
of layer~$i\!+\!1$.
The two coupling entries use different loop variables but address the
same tensor.
The dimensions that appear in both are \emph{aliased}, and their
tile sizes must satisfy cascading inequalities at every hierarchy
level~$l$:
%
\begin{align}
  X_{i+1}^{(l)} + S_{i+1}^{(l)} - 1 &\;\le\; X_i^{(l)},
    \tag{Width constraint} \\[2pt]
  Y_{i+1}^{(l)} + R_{i+1}^{(l)} - 1 &\;\le\; Y_i^{(l)},
    \tag{Height constraint} \\[2pt]
  C_{i+1}^{(l)} &\;\le\; Z_i^{(l)}.
    \tag{Channel constraint}
\end{align}

\paragraph{2.\ Decoupling.}
Per-layer weight dimensions $(Z_i, C_i, R_i, S_i)$ are orthogonal
across layers: the tiling of $R_3$ does not affect layer~5's data
reuse.
Decoupled dimensions are free to iterate at any granularity,
independent of the aliasing constraints.

\paragraph{3.\ Depth-first scheduling heuristic.}
The width, height, and channel constraints above are enforced
\emph{at every hierarchy level} for every inter-layer connection.
Every candidate mapping move triggers the heuristic validator; moves
that violate the constraints are rolled back immediately, pruning
infeasible regions of the $10^{463}$-point map space before cost
evaluation.

\paragraph{4.\ Per-layer cost model.}
The single-layer cost model is generalised along three axes:

\emph{Memory operations.}\;
Reads and writes are computed per layer per level, with per-layer
dimension filtering and bypass propagation:
%
\begin{equation}
  \text{WMOPs}_i \;=\;
    \sum_{l \in \text{levels}}
      e_l \cdot \bigl(R_l^{(i)} + W_l^{(i)}\bigr),
  \tag{Per-layer energy}
\end{equation}
%
where $e_l$ is the per-access energy at level~$l$.

\emph{Latency with stall detection.}\;
At each level~$l$, the ideal bandwidth for each layer's traffic is
compared to the physical bandwidth $B_l$.
If the ideal exceeds the physical, the level is memory-bound and the
latency is stretched:
%
\begin{equation}
  \Lambda_l^{(i)} \;=\;
    \max\!\bigl(
      \Lambda_{l,\text{rd}}^{(i)},\;
      \Lambda_{l,\text{wr}}^{(i)}
    \bigr),
  \qquad
  \Lambda \;=\;
    \max_{l}\;\sum_{i=0}^{F-1} \Lambda_l^{(i)}.
  \tag{Latency}
\end{equation}

\emph{EDP.}\;
The energy-delay product $\text{EDP} = \text{WMOPs} \times \Lambda$
provides a single-number efficiency metric that captures the trade-off
between energy and throughput.


% ====================================================================
%  §5  VALIDATION
% ====================================================================
\section*{Validation}

The extended framework is validated by reproducing the depth-first
execution of an 11-layer FSRCNN workload on the DepFiN accelerator
($16 \times 128$~PEs, 1056\,KB FMEM, 524\,KB WMEM, LPDDR4 DRAM at
18\,B/cycle read + 18\,B/cycle write, $\sim$930\,MHz clock).
The key validation results are:

\begin{enumerate}
  \item \textbf{Ideal latency}: the model predicts 37.94\,M cycles,
        matching the compute-bound reference
        ($\text{Total MACs} / \text{PEs}$) exactly.
  \item \textbf{DRAM bypass}: zero intermediate DRAM traffic is
        produced, reproducing DepFiN's depth-first fusion guarantee.
  \item \textbf{On-chip stalls}: only 24 stall cycles out of
        37.94\,M arise from the FMEM, and these occur exclusively at
        the boundary layers (L0 and L10).
        At the regime layers (L1--L9), the FMEM sustains the PE
        throughput with \emph{zero} stall cycles.
  \item \textbf{Weight Memory}: zero stall cycles, confirming that
        the 16\,B/cycle bandwidth is sufficient for weight delivery.
\end{enumerate}

\noindent
These results establish confidence in the analytical cost model for
the design-space exploration that follows.


% ====================================================================
%  §6  EXPERIMENTAL RESULTS
% ====================================================================
\section*{Experimental Results}

The exploration covers two spatial architectures (DepFiN-like,
Eyeriss-v1-like), four CNN workloads (FSRCNN, MC-CNN --- both
activation-dominant; VGG16, ResNet18 --- both weight-dominant), and
three execution modes: \emph{full fusion} (all layers fused),
\emph{partial fusion} (pairwise), and \emph{single-layer}.

\paragraph{DRAM traffic reduction.}
Full fusion reduces total DRAM traffic by 94.8\% (FSRCNN),
85.3\% (MC-CNN), 59.9\% (VGG16), and 25.1\% (ResNet18),
confirming that fusion eliminates the dominant off-chip transfers
for activation-heavy workloads.

\paragraph{Latency.}
Fusion delivers 42--48\% latency savings for the
activation-dominant workloads on DepFiN and 16--46\% on Eyeriss.
For the weight-dominant workloads, DepFiN can be \emph{slower}
under full fusion (up to 30\% latency increase) because the very
large fused weight memory requires additional scheduling passes.
Critically, partial (pairwise) fusion captures 87--103\% of the
full-fusion latency benefit, indicating that nearly all the saving
is realised by fusing just two consecutive layers.

\paragraph{Energy.}
Full fusion is the most energy-expensive mode in every
architecture--workload combination: the penalty ranges from
$1.3\times$ (FSRCNN, DepFiN) to $36.9\times$ (VGG16, Eyeriss).
The cause is the enlarged on-chip memory: per-access energy scales
with buffer capacity, and the Eyeriss architecture --- which
replicates the per-PE weight register across all 16\,384 PEs ---
suffers the most.

\paragraph{Energy-Delay Product.}
Full fusion wins on EDP in only 1 of 8
architecture--workload combinations (FSRCNN on DepFiN, where the
$47.6\%$ latency saving outweighs the $1.3\times$ energy penalty).
In all other cases, single-layer execution is more efficient.
Partial fusion, though requiring much less on-chip memory, does
not change this verdict: the energy overhead still dominates.

\paragraph{Key insight.}
Layer fusion is primarily a \emph{latency optimisation}, not an
energy optimisation.
Its benefit grows with the activation-to-weight traffic ratio and
shrinks as the on-chip memory overhead increases.
For throughput-critical deployments, partial fusion on a modestly sized
architecture offers the best compromise: it captures nearly all the
latency benefit of full fusion at a fraction of the memory cost.


% ====================================================================
%  §7  CONCLUSIONS
% ====================================================================
\section*{Conclusions}

This thesis demonstrates that the FactorFlow factor-based abstraction
generalises naturally to fused-layer mappings via aliasing and
decoupling, that the depth-first heuristic effectively prunes the
otherwise intractable $10^{463}$-point map space, and that the
per-layer cost model accurately reproduces the behaviour of the DepFiN
accelerator.
The experimental campaign clarifies the trade-offs of layer fusion:
DRAM traffic reduction of 25--95\%, latency savings of up to 48\% for
activation-dominant workloads, but an energy overhead that makes
single-layer execution the more efficient choice in 7 out of 8
tested configurations.
Partial (pairwise) fusion emerges as the most practical deployment
strategy, balancing latency savings against hardware cost.

The primary directions for future work include: extending the automated
mapper to explore fused mappings directly, supporting non-sequential
layer topologies (residual connections, multi-branch), formal
verification of the heuristic, and porting the tool to a compiled
language (C++ or Rust) to eliminate interpreter overhead and enable
shared-memory multithreading for large-scale exploration.

\vspace{1em}
\noindent\textbf{Keywords:} deep neural networks, spatial architectures,
layer fusion, mapping, design-space exploration, FactorFlow, DepFiN,
analytical cost modelling.
